%% IIT Mandi assignment template: problem sets, homework and tutorial sheets.
%% Fill in the course and student details, then add one \begin{problem} per question.
\documentclass[11pt,a4paper]{article}

\usepackage[a4paper,margin=25mm,headheight=28pt]{geometry}
\usepackage[T1]{fontenc}
\usepackage{lmodern}
\usepackage{amsmath,amssymb,amsthm}
\usepackage{graphicx}
\usepackage{xcolor}
\usepackage{booktabs}
\usepackage{enumitem}
\usepackage{fancyhdr}
\usepackage[hidelinks]{hyperref}

%% ---------------------------------------------------------------- details
\newcommand{\coursecode}{IC 152}
\newcommand{\coursename}{Computing and Data Science}
\newcommand{\assignment}{Assignment 1}
\newcommand{\studentname}{Your Name}
\newcommand{\rollnumber}{B25001}
\newcommand{\instructor}{Dr.\ Instructor Name}
\newcommand{\duedate}{15 August 2026}

%% ---------------------------------------------------------------- layout
\definecolor{iitmnavy}{HTML}{00237D}
\pagestyle{fancy}
\fancyhf{}
\fancyhead[L]{\small\coursecode: \assignment}
\fancyhead[R]{\small\studentname\ (\rollnumber)}
\fancyfoot[C]{\small\thepage}
\renewcommand{\headrulewidth}{0.4pt}

\newcounter{problem}
\newenvironment{problem}[1][]
  {\refstepcounter{problem}\par\medskip\noindent
   {\color{iitmnavy}\bfseries Problem \theproblem\if\relax\detokenize{#1}\relax\else\ (#1)\fi.}\ }
  {\par\medskip}
\newenvironment{solution}
  {\par\noindent\textit{Solution.}\ }
  {\hfill$\square$\par\medskip}

\setlength{\parindent}{0pt}
\setlength{\parskip}{4pt}

\begin{document}

%% ---------------------------------------------------------------- heading
\thispagestyle{plain}
\begin{minipage}[c]{0.18\textwidth}
  \includegraphics[width=\linewidth]{iitmandi-logo}
\end{minipage}\hfill
\begin{minipage}[c]{0.78\textwidth}
  \raggedleft
  {\Large\bfseries\color{iitmnavy}\coursecode: \coursename\par}
  \medskip
  {\large\assignment\par}
  \smallskip
  {\small Instructor: \instructor \quad\textbar\quad Due: \duedate\par}
\end{minipage}

\medskip
\noindent\rule{\textwidth}{0.8pt}
\noindent\textbf{Name:} \studentname \hfill \textbf{Roll number:} \rollnumber
\par\noindent\rule{\textwidth}{0.4pt}

%% ---------------------------------------------------------------- problems
\begin{problem}[10 marks]
Show that $\sum_{k=1}^{n} k = \frac{n(n+1)}{2}$ for every positive integer $n$.
\end{problem}

\begin{solution}
By induction on $n$. For $n = 1$ both sides equal $1$. If the identity holds
for $n$, then
\[
  \sum_{k=1}^{n+1} k = \frac{n(n+1)}{2} + (n+1) = \frac{(n+1)(n+2)}{2},
\]
so it holds for $n + 1$.
\end{solution}

\begin{problem}[5 marks]
Answer the following.
\begin{enumerate}[label=(\alph*)]
  \item First part of the question.
  \item Second part of the question.
\end{enumerate}
\end{problem}

\begin{solution}
\begin{enumerate}[label=(\alph*)]
  \item Answer to the first part.
  \item Answer to the second part, with a table if needed:

  \begin{tabular}{@{}lr@{}}
    \toprule
    Input & Output \\
    \midrule
    1 & 1 \\
    2 & 3 \\
    \bottomrule
  \end{tabular}
\end{enumerate}
\end{solution}

\end{document}
